\section{Finite-block risk and scaling}

We restate the setup. The test mass, dwell weight, innovation mass, and target
energy are $\lambda_j$, $\tau_j$, $q_j=(\lambda_j/\tau_j)/Z$, and
$s_j=\lambda_j\theta_j^2$, where $Z=\sum_j\lambda_j/\tau_j$. The number of
blocks in mode $j$ is $K_j\sim\mathrm{Binomial}(B,q_j)$. Conditional on
$K_j=k>0$, the coordinate average is unbiased and has variance
$\sigma^2/(\tau_jk)$; when $K_j=0$ it is zero.

\begin{theorem}[Exact risk]
For every $B,M\ge1$,
\[
\E\risk=\sum_{j>M}s_j+\sum_{j\le M}s_j(1-q_j)^B
+\sigma^2\sum_{j\le M}\frac{\lambda_j}{\tau_j}
\E\left[\frac{\1\{K_j>0\}}{K_j}\right].
\]
\end{theorem}

\begin{proof}
Condition on $K_j$. If $j>M$ the model output is zero and the contribution is
$s_j$. If $j\le M$ and $K_j=0$, the same contribution occurs with probability
$(1-q_j)^B$. If $K_j=k>0$, independence of the label noises gives prediction
variance $\lambda_j\sigma^2/(\tau_jk)$. Summing completes the proof.
\end{proof}

\begin{lemma}[Reciprocal binomial count]
Let $K\sim\mathrm{Binomial}(B,q)$ and
$h_B(q)=\E[\1\{K>0\}/K]$. Universal constants $c,C>0$ satisfy
\[
c\min\{Bq^2,B^{-1}\}\le qh_B(q)
\le C\min\{Bq^2,B^{-1}\}.
\]
\end{lemma}

\begin{proof}
Suppose first that $Bq\le1$. Since $1/K\le1$ on $K>0$,
\[
qh_B(q)\le q\Prob(K>0)\le Bq^2.
\]
For the reverse inequality,
\[
qh_B(q)\ge q\Prob(K=1)=Bq^2(1-q)^{B-1}.
\]
When $B=1$ the final factor equals one. When $B\ge2$, $q\le1/B$ and
$(1-q)^{B-1}\ge(1-1/B)^{B-1}\ge e^{-1}$, up to an immaterial universal
constant.

Now suppose $Bq\ge1$. For $K\ge1$, $K^{-1}\le2(K+1)^{-1}$, while
\[
\E\frac1{K+1}
=\frac{1-(1-q)^{B+1}}{(B+1)q}.
\]
Thus $qh_B(q)\le2/(B+1)\le2/B$. For the lower bound, Markov's inequality and
$\E K=Bq$ give $\Prob(K>2Bq)\le1/2$, whereas
$\Prob(K=0)=(1-q)^B\le e^{-Bq}\le e^{-1}$. Hence
\[
\Prob(1\le K\le2Bq)\ge1/2-e^{-1}>0,
\]
and on this event $K^{-1}\ge(2Bq)^{-1}$. Therefore $qh_B(q)\ge c/B$.
\end{proof}

\begin{lemma}[Weighted missing mass]
Let $p>1$, $b>1$, $q_j\asympconst j^{-p}$, and
$s_j\asympconst j^{-b}$. Then, uniformly in $B,M\ge2$,
\[
\sum_{j>M}s_j+\sum_{j\le M}s_j(1-q_j)^B
\asympconst M^{1-b}+B^{-(b-1)/p}.
\]
\end{lemma}

\begin{proof}
The tail satisfies $\sum_{j>M}s_j\asymp M^{1-b}$. Put $J=B^{1/p}$ and
$\beta=(b-1)/p$. For the upper bound, $(1-q_j)^B\le e^{-cBj^{-p}}$, and an
integral comparison gives
\[
\sum_{j\ge1}j^{-b}e^{-cBj^{-p}}
\le C B^{-\beta}.
\]
For completeness, split at $J$. The part $j>J$ is bounded by $CJ^{1-b}$.
For $j\le J$, compare to the integral and substitute $u=cBx^{-p}$; the result
is a constant multiple of
$B^{-\beta}\int_c^\infty u^{\beta-1}e^{-u}\,du$.

For the lower bound, the approximation tail already gives $cM^{1-b}$.
Choose a sufficiently large constant $A$. If $M\le2AJ$, this term is also at
least a constant multiple of $J^{1-b}$. If $M>2AJ$, then
$[AJ,2AJ]\subseteq[1,M]$ and $Bq_j\le1/2$ throughout that annulus. Hence
$(1-q_j)^B\ge(1-1/(2B))^B\ge c$, and summing $s_j$ over the annulus gives
$cJ^{1-b}=cB^{-\beta}$.
\end{proof}

\begin{lemma}[Variance sum]
Under the assumptions above,
\[
\sum_{j\le M}q_jh_B(q_j)
\asympconst \frac{\min\{M,B^{1/p}\}}{B}.
\]
\end{lemma}

\begin{proof}
By the reciprocal-count lemma, the summand is comparable to
$\min\{Bq_j^2,B^{-1}\}$. Put $J=B^{1/p}$. For $j\le J$, $Bq_j$ is bounded
below by a positive constant, so the summand is comparable to $B^{-1}$. For
$j>J$, it is comparable to $Bj^{-2p}$. If $M\le J$, summing gives $M/B$.
If $M>J$, the first part contributes $J/B$, while
$B\sum_{j>J}j^{-2p}\asymp BJ^{1-2p}=J/B$. Matching lower bounds follow from
the same regions.
\end{proof}

\begin{theorem}[Power-law risk]
If $\lambda_j\asymp j^{-a}$, $\tau_j\asymp j^r$, and
$s_j\asymp j^{-b}$, with $a>1$, $r\ge0$, and $b>1$, then for
$p=a+r$,
\[
\E\risk\asympconst M^{1-b}+B^{-(b-1)/p}
+\sigma^2\frac{\min\{M,B^{1/p}\}}{B}.
\]
\end{theorem}

\begin{proof}
Equation $q_j=(\lambda_j/\tau_j)/Z$ gives $q_j\asymp j^{-p}$ and
$\lambda_j/\tau_j=Zq_j$. Apply the preceding two lemmas to the exact identity.
\end{proof}
